\begin{center}\begin{tikzpicture}[font=\small,>=stealth,line join=round]
\draw[gray!40,rounded corners=3pt] (0,0) rectangle (14.2,8.3);\node at (7.1,7.700000000000001) {例5：两条曲线在1右侧，较大指数在上方};
\begin{scope}[shift={(1.589795918367347,2.243548387096774)},x=1.2244897959183672cm,y=1.4838709677419353cm]\draw[->,gray] (-0.4,0)--(4.5,0) node[right]{$x$};\draw[->,gray] (0,-0.4)--(0,2.7) node[above]{$y$};\draw[blue,thick,solid,domain=0.02:4,samples=180] plot(\x,{(\x)^(.3)});\draw[violet,thick,dashed,domain=0.02:4,samples=180] plot(\x,{(\x)^(.5)});\fill[orange!85!black] (1,1) circle(2.2pt);\draw[gray,dashed] (1,0)--(1,2.4);\node[below] at (1,0){$1$};\end{scope}\node[anchor=west,align=left,text width=5.0cm] at (8.4,4.0) {蓝：$x^{0.3}$\\[9pt]紫：$x^{0.5}$\\[9pt]仅讨论$x>0$\\[9pt]交点$(1,1)$不满足严格小于\\[9pt]解集$(1,+\infty)$};
\node[text width=12.6cm,align=center] at (7.1,.65) {0到1之间，蓝线反而在紫线上方；除以正数x的0.3次方也得到同一分界。};
\end{tikzpicture}\end{center}
